\chapter{The Tower Turns Around}

\section{The backward descent}

The first three chapters built \emph{up}. A separation became a count, a count became a tower of finite conditions climbing
toward a real, and the residue was carried up through every rung. This chapter turns the construction around. Measurement, it
argues, is not the endless building-up; it is the walk back \emph{down} the tower one has built --- the descent that reads,
off each rung, what the climb deposited there. For every operation that raises the construction a level there is an adjoint
that lowers it, and the downward operation is where the reading happens.

The geometry carries this duality in its bones. The exterior derivative $\mathrm{d}$ raises the degree of a form,
$\mathrm{d} : \Omega^k \to \Omega^{k+1}$, and it has an adjoint, the codifferential $\delta = \pm \star\,\mathrm{d}\,\star$,
that lowers it, $\delta : \Omega^{k} \to \Omega^{k-1}$. The climb up is $\mathrm{d}$; the descent is $\delta$; and the two
are joined at the hip by the inner product they are adjoint in, $\langle \mathrm{d}\alpha, \beta \rangle = \langle \alpha,
\delta\beta \rangle$. Every structure the instrument built going up with $\mathrm{d}$ --- the potential from which the
curvature came, the tower of forms --- has a reading going down with $\delta$. The Laplacian that governs the whole space,
$\Delta = \mathrm{d}\delta + \delta\mathrm{d}$, is built from both directions at once: a climb and its descent, composed. To
read the tower is to apply the downward operator, and the physics of the space lives in that composition.

There is a second face of the same duality, and it is the one the physicist meets first: the pullback. A map $f : N \to M$
pushes points forward but pulls forms \emph{back}, $f^\star : \Omega^k(M) \to \Omega^k(N)$, carrying structure down the map
against the direction of the points. This is the descent in its most useful form. The instrument builds a field on the large
space and reads it by pulling it back onto the small probe --- the worldline, the bench, the loop --- where the reading
actually happens. Measurement is a pullback: the world's field, defined up in the ambient geometry, restricted down onto the
apparatus that meets it. And the pullback commutes with the climb, $f^\star \mathrm{d} = \mathrm{d} f^\star$, so the residue
built up top survives the descent onto the probe unchanged --- which is exactly why the descent can read it.

The reason the descent, and not the climb, is where measurement lives is that building is unbounded and reading is finite. You
can always take one more exterior derivative, add one more rung, climb one more level; the climb has no natural stopping
place. The descent does: it stops at the probe, at degree zero, at the number the bench returns. Riemann's own procedure was of
this kind --- he built the metric relations up over the bare manifold and then read curvature back down as the local
obstruction to flatness~\cite{riemann1854}. The instrument inherits the same shape. It does not measure by climbing forever;
it measures by climbing to where the residue is deposited and then walking back down to read it off. The turn from up to
down is the turn from construction to measurement.

In the two verbs, the climb is a funge that accumulates and the descent is the tange that reads. Going up, the instrument
funges rung onto rung, gathering the tower; coming down, it tanges off each rung the one characteristic the probe can pull
back, $f^\star$, the contravariant restriction onto the apparatus. The adjointness $\langle \mathrm{d}\alpha, \beta\rangle =
\langle \alpha, \delta\beta\rangle$ is the statement that the funge going up and the tange coming down are the same pairing
read in two directions. So the tower turns around exactly where funge hands off to tange --- where the gathered construction
is finally selected from, rung by rung, on the way down. And the first thing the descent finds is the seam, the join between
one rung and the next, where the residue was deposited.

\section{The seam}

The descent walks down the tower rung by rung, and between every two rungs there is a join. This section is about that join
--- the seam where level $n$ hands the construction to level $n+1$, where the residue is deposited, and where the instrument
is billed. Nothing is free at a seam. Each turn from one level to the next carries a cost, the cost is the residue, and the
total bill of a measurement is the sum of its turns. The seam is small and it is everywhere, and reading it is the whole of
what the descent does.

The geometry of a seam is the transition between charts, or the composition of transports. Where two charts overlap, the
construction on one must agree with the construction on the other up to a transition function $g_{\alpha\beta}$, and the way
a connection carries a state across the overlap is by parallel transport, $P_\gamma = \mathcal{P}\exp\!\big(\!-\!\oint_\gamma
A\big)$. Compose two segments and the transports multiply; the residue accumulates in the product. Crucially, transport
around a small closed seam does not return the identity but the curvature, $P_{\square} \approx \mathbb{1} - F_{\mu\nu}\,
\mathrm{d}x^\mu \wedge \mathrm{d}x^\nu$ to leading order. So the join between rung $n$ and rung $n+1$ deposits exactly the
residue of Chapter 3: the curvature lives at the seam, and to cross the seam is to pick it up. Pass $n$ genuinely seeds pass
$n+1$ --- the state handed across is the state carried by the transport --- and what is added in the handing is the residue.

This is what it means to say the instrument is \emph{billed the turns}. The cost of a construction is not in its rungs, which
are just states, but in its joins, where the transports compose and the curvature is deposited. A path with many seams
accumulates much residue; a path with few accumulates little; and the total is the ordered product of the turns, path by
path. The physics reads this as the plain fact that a phase, an energy, an angle is accumulated along a route and depends on
the route --- that two paths between the same endpoints can return different residues if they enclose different curvature.
The bill is not levied on being somewhere; it is levied on the turns taken to get there. The instrument pays for its seams,
never for its rest.

There is a subtlety the construction is careful about, and it is the one that keeps the seam honest. The residue at a seam is
not a matter of how the labels were chosen; it is the covariant curvature, invariant under any relabeling of the charts.
Change the transition function by a gauge transformation and the transport changes by conjugation, $P \mapsto h\,P\,h^{-1}$,
but the loop's residue --- its trace, its holonomy class --- is untouched. So the bill charged at a seam is a real cost, not
an artifact of bookkeeping. You cannot relabel your way out of paying it, any more than you can gauge away the class $[F]$.
This is the seam's version of the last chapter's lesson: what is deposited at the join is the invariant, and the invariant is
what you are charged for.

In the two verbs, the seam is where a funge is paid for with a tange. Crossing from rung $n$ to rung $n+1$, the instrument
funges the two levels into one continued construction --- bags them as a single tower --- and is charged, for the bagging, a
tange: the curvature $F_{\mu\nu}$ selected out at the join, the covariant residue the crossing deposits. The bill is the
funge's price, and the price is a residue that no relabeling can waive. So the descent, reading down the tower, reads at each
seam exactly the turn's charge --- and the sum of those charges, over the whole descent, is what the loop will hand back
altered. That alteration, gathered from all the seams, is the subject of the last section.

\section{What came back changed}

The descent has walked down the tower and read, at each seam, the residue the turn deposited. This closing section gathers
those readings into the one idea the whole part exists to establish: measurement is a loop, and what a measurement returns is
not the tower but the \emph{alteration} --- the difference between what set out and what came back. The instrument does not
measure by building forever and reading the top. It measures by closing a circuit and asking what returned changed. The
change is the reading; everything else is scaffolding.

Assemble the seams into a loop and the accumulated residue is the holonomy, $P_\gamma = \mathcal{P}\exp\!\big(\!-\!\oint_
\gamma A\big)$, and its departure from the identity is precisely what the loop hands back altered. A state set out around the
circuit; the transport carried it seam by seam; and it returned rotated, phased, shifted by the enclosed curvature. That
returned-altered state is the measurement. In the flat case the loop closes to the identity and nothing is measured --- there
was no residue to read. In the curved case the loop returns a rotation, and the size of that rotation is the number the bench
records. So a measurement is, structurally, a comparison of the sent and the returned across a closed circuit, and its output
is their difference.

This reframes the whole activity of the instrument, and the reframing is the payoff of the part. The naive picture is that
measurement builds a taller and taller tower --- more precision, more rungs, a real approached ever more closely --- and
reads the answer at the top. The honest picture is that measurement closes a loop and reads the alteration at the bottom. The
first picture has no stopping place, because the tower can always grow; the second has a natural one, because the loop
closes. And the second is what the physics actually does: an interferometer sends a state around two arms and reads the phase
between them; a gyroscope carries an axis around a circuit and reads the angle it returned rotated; a bench measures a
coupling by comparing a probe before and after it has been around the field. In each case the measurement is the residue of a
closed loop, not the height of an open tower.

The construction supplies the loop and its holonomy; the identification of the returned-altered state with a specific
physical reading --- a phase, an angle, a coupling --- is the gauge's, and marked as such. What the instrument derives is
that a closed circuit through curvature returns altered by an invariant residue. What the physics reads is which observable
that residue is, at each bench. The book keeps them apart, as always, so that the loop's alteration is never confused with a
magnitude derived from first principles. The loop returns changed; how much it changed, and whether that number can ever be
pinned rather than bracketed, is the question the next chapters open.

In the two verbs, the loop is a funge the instrument sends around and a tange it reads on return. The circuit funges the
whole path into one gathered transport --- bags every seam into a single holonomy --- and the return tanges off the one
characteristic that changed, the rotation the curvature deposited. Measurement is that pair: funge the loop, tange the
alteration. What came back the same is the funge that closed; what came back changed is the tange the residue forced. So Part
II ends with measurement re-understood as the reading of a loop's alteration --- and the very next chapter takes the place
where the alteration first refuses to be a clean rotation, where the loop's failure to close becomes, for the first time, a
number.
